\chapter[پیاده‌سازی سرویس‌دهنده]{طراحی و پیاده‌سازی سرویس‌دهنده}

این فصل کار انجام‌شده توسط \textbf{مهیار عامری} را شرح می‌دهد: خدمت پشتیبان با \textenglish{Django}، لایهٔ \textenglish{REST} با \textenglish{Django REST Framework}، و منطق دامنه روی \textenglish{ORM} جنگو.

\section{انتخاب فناوری}

جنگو برای این پروژه به سه دلیل انتخاب شد. نخست، \textenglish{ORM} بالغ آن مدل‌های حمل‌ونقل را بدون \textenglish{SQL} دستی نگه می‌دارد. دوم، \textenglish{DRF} سریالایزر، مجوز، صفحه‌بندی و مستندسازی \textenglish{OpenAPI} را یکجا می‌دهد. سوم، استقرار روی سکوی لیارا برای کاربرد آموزشی با کمترین پیکربندی عملی است. پایگاه پیش‌فرض \textenglish{SQLite3} با تنظیمات استاندارد جنگو برای حجم فعلی داده کافی است و مهاجرت به موتور دیگر تنها با تغییر تنظیمات اتصال انجام می‌شود.

\section{سازمان پروژهٔ جنگو}

پروژه از الگوی چنداپلیکیشنی جنگو پیروی می‌کند. اپلیکیشن حساب‌ها کاربر سفارشی (یا پروفایل متصل به کاربر) را با شماره موبایل به‌عنوان شناسهٔ ورود مدیریت می‌کند. اپلیکیشن حمل، مدل بار و تخصیص را دارد. اپلیکیشن مرجع، کالا، ماشین و استان را جدا نگه می‌دارد تا فهرست‌های فرم در کلاینت پایدار بماند.

تنظیمات اصلی:

\begin{itemize}
  \item \en{INSTALLED\_APPS} شامل \en{rest\_framework} و بستهٔ ژتون ساده (\textenglish{Simple JWT}) است.
  \item احراز هویت پیش‌فرض \textenglish{DRF} روی \en{JWTAuthentication} تنظیم شده است.
  \item صفحه‌بندی از نوع شماره صفحه است؛ پاسخ شامل \en{results}، \en{current\_page}، \en{total\_pages} و \en{next} است تا کلاینت بتواند همهٔ مأموریت‌ها را حلقه کند.
  \item \en{LANGUAGE\_CODE} و منطقه متناسب با ایران است؛ خودِ کد وضعیت بار انگلیسی می‌ماند تا قرارداد پایدار بماند.
\end{itemize}

\section{احراز هویت و صدور ژتون}

ورود \en{POST /api/accounts/login} شماره و رمز را می‌گیرد، کاربر را با \en{authenticate} جنگو بررسی می‌کند و در صورت موفقیت دو چیز برمی‌گرداند: شیء کاربر (شامل \en{is\_driver} و \en{is\_operator}) و ژتون‌ها. ثبت‌نام \en{POST /api/accounts/register/} کاربر را می‌سازد، رمز را با \en{set\_password} درهم می‌کند و برای راننده \en{machine\_id} و \en{ostan\_id} را روی پروفایل می‌نویسد. ثبت‌نام ژتون نمی‌دهد تا ساخت حساب از شروع نشست جدا بماند.

پروفایل \en{GET/PUT /api/accounts/profile/} فیلدهای نام، کد ملی، و برای راننده پلاک (\en{pelak})، مدل، ماشین و استان را به‌روز می‌کند. تغییر رمز \en{old\_password} را با \en{check\_password} می‌سنجد و سپس رمز جدید را ذخیره می‌کند. خروج در سمت سرور می‌تواند ژتون تازه‌سازی را در فهرست سیاه بگذارد؛ کلاینت در هر حال ذخیره‌ساز محلی را پاک می‌کند.

مجوز نقش با کلاس‌های \textenglish{DRF} پیاده شده است: ویوهای \en{/api/driver/...} فقط برای \en{is\_driver} و ویوهای \en{/api/operator/...} فقط برای \en{is\_operator} باز هستند. کاربر راننده نمی‌تواند بار متصدی دیگری را حذف کند؛ queryset متصدی با \en{request.user} فیلتر می‌شود.

\section{منطق بار و تخصیص}

ایجاد بار فقط از مسیر متصدی مجاز است. سریالایزر فیلدهای عنوان، توضیحات، قیمت، وزن، کالا، ماشین، استان‌ها و نشانی‌ها را اعتبارسنجی می‌کند. مختصات جغرافیایی اختیاری است؛ اگر کلاینت نفرستد، سرور می‌تواند بعداً از نشانی پر کند یا خالی بگذارد. وضعیت اولیه \en{open} است.

فهرست راننده فقط بارهای \en{open} را برمی‌گرداند و فیلترهای \en{ostan\_mabda}، \en{ostan\_maghsad}، \en{price\_min} و \en{price\_max} را روی queryset اعمال می‌کند. پذیرش بار در یک تراکنش انجام می‌شود: قفل ردیف بار، بررسی وضعیت باز بودن، بررسی سازگاری ماشین راننده با ماشین بار، ساخت رکورد تخصیص، و تغییر وضعیت به \en{assign}. اگر بار پیش‌تر گرفته شده باشد، پاسخ خطا برمی‌گردد.

تکمیل کار دو نقطهٔ پایانی جدا دارد. هر کدام فقط پرچم مربوط به نقش را \en{True} می‌کند. وقتی هر دو پرچم درست باشند، وضعیت بار \en{done} می‌شود. این کار با \en{update} روی مدل تخصیص و در صورت لزوم \en{save} روی بار، داخل \en{transaction.atomic} انجام می‌شود تا حالت میانی دیده نشود.

\section{سریالایزرها}

\textenglish{DRF} برای هر مورد یک سریالایزر خواندن و یک سریالایزر نوشتن دارد. در خواندن بار، نام‌های نمایشی \en{product\_name}، \en{machine\_name}، \en{ostan\_mabda\_name}، \en{operator\_name} و در صورت وجود \en{assign\_info} به صورت تو در تو اضافه می‌شوند تا کلاینت مجبور به درخواست اضافه نباشد. در نوشتن، فقط شناسهٔ کالا و ماشین و استان پذیرفته می‌شود. وزن به صورت عدد اعشاری ذخیره می‌شود؛ برای سازگاری با نسخه‌های قدیمی، توضیحات ممکن است پسوند وزن هم داشته باشد.

\section{خطاها و پیام‌ها}

پاسخ خطا به صورت \textenglish{JSON} با کلید متعارف (\en{detail} یا \en{error} یا خطاهای فیلدی سریالایزر) برمی‌گردد. کلاینت این پیام‌ها را با \en{ApiResponse} به متن قابل نمایش تبدیل می‌کند. کدهای ۴۰۱ برای ژتون نامعتبر، ۴۰۳ برای نقش نادرست، ۴۰۴ برای بار ناموجود و ۴۰۰ برای اعتبارسنجی فرم استفاده می‌شوند.

\section{استقرار سرور}
\enlargethispage{3.6\baselineskip}

نسخهٔ در حال بهره‌برداری آموزشی روی لیارا با دامنهٔ زیر در دسترس است:
\siteline{https://tran-develoop.liara.run}
جنگو در این محیط با \textenglish{WSGI/ASGI} پشت پیشکار اجرا می‌شود. فایل پایگاه \textenglish{SQLite} روی دیسک پایدار سرویس قرار می‌گیرد. کلیدهای نشان اگر از سرور پروکسی شوند در متغیر محیطی می‌مانند و به کلاینت وب داده نمی‌شوند.
